\subsection{Cyclic prime outer actions and the full \texorpdfstring{$\operatorname{PSL}_2(32)$}{PSL2(32)} chief step}

\begin{statement}{D.45}{Proved}{All multiplier-one peripheral fibres across split cyclic-prime outer chief extensions}
Let \(S\) be a finite nonabelian simple group, and fix an
actual isomorphism
\[
 \operatorname{Aut}(S)=S\rtimes\langle t\rangle,\qquad
 |t|=r,\quad
 \operatorname{Out}(S)=\langle\bar t\rangle\simeq C_r,
 \quad r\ \text{prime},
 \tag{D.45.1}
\]
where \(S\) is identified with its inner automorphisms.
Let
\[
 1\longrightarrow A=S^\Omega\longrightarrow E
       \xrightarrow{\pi}D\longrightarrow1
 \tag{D.45.2}
\]
be a finite chief extension, and let
\(q_E:\widehat F(x_1,x_2)\twoheadrightarrow E\) be an
actual epimorphism, with \(x_3=(x_1x_2)^{-1}\).
Put \(X_i=q_E(x_i)\), \(q_0=\pi q_E\), and \(d_i=q_0(x_i)\).
For every epimorphism \(q_1:\widehat F\twoheadrightarrow D\)
and every specified \(h_i\in D\) satisfying
\[
                    q_1(x_i)=h_id_ih_i^{-1},
 \tag{D.45.3}
\]
choose arbitrary lifts \(H_i\in E\) of \(h_i\).
There is an epimorphism \(f:\widehat F\twoheadrightarrow E\)
such that
\[
 \pi f=q_1,\qquad
 f(x_i)=a_iH_iX_iH_i^{-1}a_i^{-1}
                 \quad(a_i\in A).
 \tag{D.45.4}
\]
The number of these literal epimorphisms is exactly the
number for \((q_1,h_1,h_2,h_3)=(q_0,1,1,1)\).
No restriction is imposed on the finite base \(D\),
its factor stabilizers, or its fixed peripheral points.

In particular \textup{(D.45.1)} holds for
\[
 S=\operatorname{PSL}_2(32)=\operatorname{SL}_2(\mathbf F_{32}),
 \qquad
 t:(a_{ij})\longmapsto(a_{ij}^{\,2}),
 \qquad |\operatorname{Aut}(S)|=163680.
 \tag{D.45.5}
\]
Every D.37 automorphism \(\beta_Q\) of every specified
characteristic target \(P_n\) therefore lifts through
every actual \(\operatorname{PSL}_2(32)\)-chief extension
above that target, with its three stated base conjugators.
The maximal pro-five quotient of this extension is exactly
the given \(P_n\), with the action \(\beta_Q\).
The same conclusion retains a prescribed epimorphism into
any additional finite group over \(D\) whose kernel over
\(D\) is solvable, through their actual fibre product.
\end{statement}

\begin{proof}
We first prove the surface maps which transport the outer
phases, and then retain all the finite extension and source
maps used to count the full fibre.

\paragraph{An exact positive boundary push.}
Use a surface presentation
\[
 \Pi=\left\langle a_1,b_1,\ldots,a_g,b_g,c_1,\ldots,c_b
          \,\middle|\,
       \prod_{j=1}^g[a_j,b_j]\prod_{i=1}^b c_i=1
        \right\rangle,\qquad b\geq2.
 \tag{D.45.6}
\]
The finite-index pants covers below have \(b\geq3\).
One boundary may be eliminated to obtain a free basis;
all substitutions below preserve the displayed relation
before or after that elimination.

For one handle and an adjacent boundary \(c\), the
following are actual free-group substitutions:
\[
 \begin{split}
 P(a)&=c^{-1}ac,\\
 P(b)&=c^{-1}a^{-1}cab c,\\
 P(c)&=c^{-1}a^{-1}cac .
 \end{split}
 \tag{D.45.7}
\]
They fix \([a,b]c\) literally, and fix all the other
generators. Their inverse can be obtained from the
following explicit Nielsen changes. Put
\[
 u=ac,\qquad v=c^{-1}bc,\qquad z=vcv^{-1}.
 \tag{D.45.8}
\]
The inverse expressions are
\[
 a=uv^{-1}z^{-1}v,\qquad
 b=v^{-1}zvz^{-1}v,\qquad c=v^{-1}zv,
 \tag{D.45.9}
\]
and \([u,v]z=[a,b]c\). The twist \(u\mapsto u\),
\(v\mapsto vu\), \(z\mapsto z\), transported through
\textup{(D.45.8)--(D.45.9)}, composed with the inverse
of \(a\mapsto a,\ b\mapsto ba,\ c\mapsto c\), is
exactly \textup{(D.45.7)}. Reversing these operations
gives its inverse. This also realizes it by the two
boundary-fixed Dehn twists used in D.41.
Its abelianization sends
\[
            a\longmapsto a,\qquad b\longmapsto b+c,
                         \qquad c\longmapsto c.
 \tag{D.45.10}
\]
Its positive boundary conjugator is \(c^{-1}a^{-1}\).

For two adjacent boundaries \(c,d\), the Hurwitz map
\[
       J(c)=cdc^{-1},\quad J(d)=c,\qquad
       J^{-1}(c)=d,\quad J^{-1}(d)=d^{-1}cd
 \tag{D.45.11}
\]
fixes their product. Its square is
\[
 J^2(c)=(cd)c(cd)^{-1},\qquad
 J^2(d)=cdc^{-1}.
 \tag{D.45.12}
\]
These equations retain both positive orientations and
both conjugating paths.

\paragraph{Independent phases with a prescribed character.}
Write \(\widetilde S=S\rtimes\langle t\rangle\), with
outer projection \(\epsilon:\widetilde S\to\mathbf F_r\).
Let \(\chi:\Pi\to\mathbf F_r\) be nonzero.
For each boundary \(c_i\), prescribe an \(S\)-conjugacy
invariant subset \(T_i\subset\epsilon^{-1}(\chi(c_i))\).
The number of epimorphisms
\(\varphi:\widehat\Pi\twoheadrightarrow\widetilde S\)
with \(\epsilon\varphi=\chi\) and
\(\varphi(c_i)\in T_i\) is unchanged upon conjugating
each \(T_i\) independently by any power of \(t\).
We prove this by actual automorphisms of the source.

First suppose \(\chi(c)=0\). If a handle has nonzero
character on \(a\), then \textup{(D.45.7)} preserves
\(\chi\) by \textup{(D.45.10)} and changes the
boundary conjugator phase by \(-\chi(a)\).
If only the handle value at \(b\) is nonzero, use
the actual basis change \(a\mapsto ab,\ b\mapsto b\),
whose inverse is \(a\mapsto ab^{-1}\).
It fixes \([a,b]\); conjugating the displayed push
by this change gives phase \(-\chi(b)\).
Each nonzero phase generates \(\mathbf F_r\), so
powers of these maps give every required phase.

If the character is supplied instead by another
boundary \(d\), then \textup{(D.45.12)} preserves
\(\chi\). Its conjugator phases are \(\chi(d)\)
at \(c\) and \(\chi(c)=0\) at \(d\); all other
boundaries are fixed. Again the first phase is a
generator of \(\mathbf F_r\).
An arbitrary boundary can be moved to the indicated
position by the positive Hurwitz maps
\textup{(D.45.11)}.
Handle blocks can also be interchanged by an actual
positive Hurwitz map on their commutators: for
\(A=[a,b]\), \(B=[u,v]\), send
\[
 (a,b,u,v)\longmapsto
       (A uA^{-1},A vA^{-1},a,b).
 \tag{D.45.13}
\]
Its commutator blocks become \((ABA^{-1},A)\),
so their product is fixed. The inverse uses the
inverse Hurwitz map and conjugates both entries of
the affected handle pair by the same word.
These moves put the selected handle last and \(c\)
first among the boundaries, where
\textup{(D.45.7)} applies. Pulling the push back by
the actual inverse moves proves the assertion in
the original presentation.

Here is also the precise phase comparison under
these changes of paths. For a character-preserving
automorphism sending \(c_i\) to \(u_ic_iu_i^{-1}\),
the phase is \(\chi(u_i)\).
At a zero-character boundary it is independent of
the conjugator choice: the boundary is primitive,
and its free-group centralizer is
\(\langle c_i\rangle\), of character zero.
The phase at a fixed boundary adds under composition
of character-preserving maps.
Upon conjugation by one of the coordinate changes
above, the character is transported with that
change, and the phases of the change and its inverse
cancel. Thus the pulled-back map has the stated
phase only at the chosen zero-character boundary.
No uncontrolled based-path phase is introduced.

For a boundary with \(j=\chi(c_i)\ne0\), every
\(S\)-class in its outer coset is already invariant
under conjugation by every \(t^k\).
Indeed, for \(a\in\epsilon^{-1}(j)\), choose
\(m\) with \(mj=k\) in \(\mathbf F_r\).
Then \(t^ka^{-m}\in S\), and
\[
       t^k a t^{-k}
        =(t^ka^{-m})\,a\,(t^ka^{-m})^{-1}.
 \tag{D.45.14}
\]
This proves the assertion element by element.

Finally, if a source map has boundary conjugator
word \(u_i\) of phase \(k\), then
\(\varphi(u_i)=s_i t^k\) for an actual \(s_i\in S\).
Hence precomposition changes membership in \(T_i\)
precisely to membership in \(t^kT_it^{-k}\);
its \(S\)-factor has no effect on that set.
Precomposition by the displayed inverse gives
the reverse map. Every such map is an automorphism
of the same surface group, so surjectivity is
preserved. This proves the claimed equality of
epimorphism counts, including empty sets, by
explicit bijections.

\paragraph{Matching the actual surface characters.}
Two finite covers of the pants with the same index
and matched oriented boundaries have the same
genus. The adjacent maps \textup{(D.45.11)} realize
the prescribed boundary permutation.
Suppose two characters to \(\mathbf F_r\) have the
same values on those boundaries, and are either
both zero or both nonzero.
There is an actual surface isomorphism matching
the characters and the oriented boundary classes.
If all boundary values vanish, the nonzero
characters are nonzero covectors on the handle
symplectic space. The elementary Dehn twists of
D.41 generate the symplectic group modulo \(r\),
which is transitive on its nonzero covectors.
Their integral elementary matrices and inverses,
including the two-by-two diagonal factorization
in D.41, realize the needed finite-field map by a
finite composition of these actual homeomorphisms.

If a boundary \(d\) has character value \(e\ne0\),
the bounding pair with homology classes \(a\) and
\(a+d\) acts by
\[
                  v\longmapsto v+\langle v,a\rangle d.
 \tag{D.45.15}
\]
It is the product of the two indicated Dehn twists;
their disjoint simple curves bound that boundary
and an annulus with a handle attached.
Its \(k\)-th power changes the character by
\(ke\langle-,a\rangle\).
Since \(e\) is a unit, the handle basis curves and
their powers give every desired translation of
the handle character. In genus zero the boundary
values already determine the character.
All these maps are based outside their supports.
They induce inverse maps on the discrete surface
groups and on their profinite completions.
The extra boundary conjugators of such an
isomorphism have definite character phases; the
preceding independent-phase bijections handle
every one of them.

\paragraph{The actual twisted chief extension, with its signs.}
Let \(D_*\) be a factor stabilizer and identify
\(\Omega=D_*\backslash D\) with its right action.
Choose representatives \(t_\omega\in D\), with
\(t_{\omega_0}=1\). Centerlessness gives the actual
isomorphism
\[
 E\longrightarrow
       \operatorname{Aut}(A)
           \mathbin{\times_{\operatorname{Out}(A)}}D,
       \qquad e\longmapsto(\operatorname{Ad}_e,\pi(e)).
 \tag{D.45.16}
\]
Its kernel is \(Z(A)=1\); for a required pair,
choose a lift of its \(D\)-coordinate and correct
the difference by the unique inner automorphism
of \(A\). This proves surjectivity as well.

Write the outer action as the cocycle
\(b_d(\omega)\in\mathbf F_r\), so
\[
 b_{de}(\omega)=b_d(\omega)+b_e(\omega d).
 \tag{D.45.17}
\]
Set
\[
 \begin{split}
 \epsilon_*(c)&=b_c(\omega_0)\quad(c\in D_*),\\
 s_\omega&=b_{t_\omega}(\omega_0),\\
 k(\omega,d)&=t_\omega d t_{\omega d}^{-1}.
 \end{split}
 \tag{D.45.18}
\]
Then \(\epsilon_*:D_*\to\mathbf F_r\) is a
homomorphism, and the cocycle identity gives
\[
           b_d(\omega)=-s_\omega+
                      \epsilon_*(k(\omega,d))+s_{\omega d}.
 \tag{D.45.19}
\]
The two signs in this formula are retained.
Conjugating the wreath coordinates by
\((t^{s_\omega})_\omega\), namely
\[
          (z,d)\longmapsto
       (t^{s_\omega}z(\omega)t^{-s_{\omega d}},d),
 \tag{D.45.20}
\]
identifies the given extension with
\[
 \begin{split}
 E_*={}&\{(z,d)\in\widetilde S^\Omega\rtimes D:
          \epsilon(z(\omega))=\epsilon_*(k(\omega,d))\},\\
 (z,d)(w,e)={}&(z(\omega)w(\omega d),de).
 \end{split}
 \tag{D.45.21}
\]
This is an actual isomorphism, carrying the specified
\(X_i,H_i\) with it.

For either source \(q=q_j\), put
\(K_j=\ker q_j\) and \(H_{j,C}=q_j^{-1}(C)\).
Choose actual source words \(T_\omega\) with
\(q(T_\omega)=t_\omega\) and \(T_{\omega_0}=1\).
The twisted Shapiro formula is
\[
 \begin{split}
 \varphi:H_{j,D_*}&\longrightarrow\widetilde S,
                   \qquad \epsilon\varphi=\epsilon_*q_j,\\
 z_\varphi(g)(\omega)
    &=\varphi(T_\omega gT_{\omega q_j(g)}^{-1}),\\
 f_\varphi(g)&=(z_\varphi(g),q_j(g)).
 \end{split}
 \tag{D.45.22}
\]
The middle representatives cancel in a product,
proving its homomorphism identity and its exact
outer coordinates.
Conversely, for a cocycle \(z\) evaluate it at
\(\omega_0\) on \(H_{j,D_*}\).
The vector
\(v(\omega)=z(T_\omega)(\omega_0)\) belongs to
\(S^\Omega\) and has \(v(\omega_0)=1\), since its
outer transition at \(t_\omega\) is the identity.
It satisfies
\[
          z_\varphi(g)(\omega)
             =v(\omega)z(g)(\omega)v(\omega q_j(g))^{-1}.
 \tag{D.45.23}
\]
This proves the normalization and its inverse.
There are exactly \(|S|^{|\Omega|-1}\) literal
cocycles for each normalized surface map.

\paragraph{Every intermediate character and the prescribed conjugators.}
For \(D_*\leq C\leq D\), retain every homomorphism
\(\epsilon_C:C\to\mathbf F_r\) extending
\(\epsilon_*\); there may be none.
Choose representatives \(R_\nu\) for \(C\backslash D\),
with \(R_{C}=1\). Define
\[
 \begin{split}
 k_C(\nu,d)&=R_\nu dR_{\nu d}^{-1},\\
 b_C(d,\nu)&=\epsilon_C(k_C(\nu,d)),\\
 p(\omega)&=Ct_\omega,\qquad
 g_\omega=\epsilon_C(t_\omega R_{p(\omega)}^{-1}).
 \end{split}
 \tag{D.45.24}
\]
Let \(E_C\) be \textup{(D.45.21)} formed with
\(C\backslash D,\epsilon_C,k_C\).
The actual injection over \(D\) is
\[
 j_C:E_C\hookrightarrow E_*,
 \qquad
 j_C(z,d)(\omega)
       =t^{g_\omega}z(p(\omega))t^{-g_{\omega d}} .
 \tag{D.45.25}
\]
The adjacent powers cancel under multiplication.
Reading any coordinate above each \(\nu\) proves
injectivity. Its outer-coordinate equation is
\[
 \epsilon_*(k(\omega,d))
       =g_\omega+b_C(d,p(\omega))-g_{\omega d},
 \tag{D.45.26}
\]
which follows by applying \(\epsilon_C\) to the
literal factorization in \(C\).
Also \(g_{\omega_0}=0\).
Thus a surface homomorphism on \(H_{j,C}\),
followed by its Shapiro map and \(j_C\), restricts
at the base coordinate to its very same map on
\(H_{j,D_*}\).

For fixed \(i\), two elements over \(q_j(x_i)\)
are conjugate by \(A\) precisely when each of their
permutation-cycle products is conjugate by \(S\).
The successive conjugator equations along a cycle
give this equivalence: once its first conjugator
is chosen in \(S\), all others are determined;
their outer exponents vanish because the two
prescribed component patterns agree. The last
equation is exactly the cycle-product condition.

This supplies an \(S\)-conjugacy invariant subset
\(T_{j,C,\delta}\subset\widetilde S\) for every
boundary cycle \(\delta\) on \(C\backslash D\).
Explicitly choose the initial coordinate \(\nu\)
of \(\delta\), and initial coordinates \(\omega\)
of all cycles \(\gamma\) above it with
\(p(\omega)=\nu\).
If their relative degrees are \(m_\gamma\),
intersect the conditions
\[
        t^{g_\omega}a^{m_\gamma}t^{-g_\omega}
                   \in\mathcal C_{j,\gamma},
       \qquad
       \epsilon(a)=\epsilon_C(q_j(c_\delta)).
 \tag{D.45.27}
\]
Here \(\mathcal C_{0,\gamma}\) is the actual
\(S\)-class of the cycle product of \(X_i\), and
\(\mathcal C_{1,\gamma}\) that of
\(H_iX_iH_i^{-1}\). All roots and all empty
intersections are retained.

The exact boundary bijection is
\[
 C d\langle q_1(x_i)\rangle
       \longmapsto C d h_i\langle d_i\rangle.
 \tag{D.45.28}
\]
It preserves cycle lengths and every intermediate
boundary degree by \textup{(D.45.3)}.
For a new initial coordinate \(\nu'\), take
\(\nu'h_i\) for its old match.
In \(E_C\) the canonical element over \(h_i\)
has components
\[
                  L_i(\nu)=t^{\,b_C(h_i,\nu)}.
 \tag{D.45.29}
\]
Its image under \(j_C\) and the specified \(H_i\)
have the same projection to \(D\); their ratio
lies in \(A\). Conjugation by \(L_i\) therefore
transports the entire old local condition to the
new one, including the intersection in
\textup{(D.45.27)}.
On the cycle product the conjugator is the actual
component \(L_i(\nu')\), giving
\[
       T_{1,C,\delta'}
        =t^{\,b_C(h_i,\nu')}
              T_{0,C,\delta}\,t^{-b_C(h_i,\nu')}.
 \tag{D.45.30}
\]
The inverse conjugation proves the reverse
inclusion. This argument does not require
\(X_i\in j_C(E_C)\).
In particular the matching boundary characters
are equal, not merely proportional.

The two actual surface characters
\(\chi_j=\epsilon_Cq_j\) on \(H_{j,C}\) are both
zero when \(\epsilon_C=0\), and both nonzero
otherwise, since \(q_j(H_{j,C})=C\).
Use the character-matching isomorphism proved
above. When \(\epsilon_C=0\), its boundary
conjugators have zero outer phase, and
\textup{(D.45.30)} has phase zero as well.
When \(\epsilon_C\ne0\), the independent-phase
automorphisms remove every phase in
\textup{(D.45.30)} and every additional phase of
the matching isomorphism. They preserve
surjectivity.
Consequently the numbers of epimorphisms
\[
 H_{j,C}\twoheadrightarrow
                S\rtimes\operatorname{im}(\epsilon_C)
 \tag{D.45.31}
\]
with the exact outer character and all boundary
sets \textup{(D.45.27)} are equal for \(j=0,1\).

\paragraph{Subtracting exactly the maps factoring through the base.}
For an epimorphism in \textup{(D.45.31)},
\(\varphi(K_j)\) lies in \(S\) and is normal in
the target, hence is either \(1\) or \(S\).
We retain only those with image \(S\).
The subtracted maps are exactly the actual
epimorphisms
\[
       \Phi:C\twoheadrightarrow
                  S\rtimes\operatorname{im}(\epsilon_C),
       \qquad \epsilon\Phi=\epsilon_C,
       \qquad \varphi=\Phi q_j|_{H_{j,C}}.
 \tag{D.45.32}
\]
There is a bijection of the subtracted sets
using the same \(\Phi\) on both sides.
Indeed its Shapiro section is the homomorphism
\[
 s_\Phi:D\longrightarrow E_C,\qquad
 s_\Phi(d)(\nu)=\Phi(R_\nu dR_{\nu d}^{-1}),
 \tag{D.45.33}
\]
with base coordinate \(d\).
It is independent of \(q_j\).
The two normalized source maps are
\(j_Cs_\Phi q_j\).
At a peripheral, \textup{(D.45.3)} gives
\[
 j_Cs_\Phi(q_1(x_i))
   =j_Cs_\Phi(h_i)\,
             j_Cs_\Phi(d_i)\,
                         j_Cs_\Phi(h_i)^{-1}.
 \tag{D.45.34}
\]
The first factor and \(H_i\) differ by an element
of \(A\).
Thus this value is \(A\)-conjugate to
\(H_iX_iH_i^{-1}\) exactly when
\(j_Cs_\Phi(d_i)\) is \(A\)-conjugate to \(X_i\).
This proves both directions of the desired
bijection for the very same \(\Phi\).
Subtracting these equal finite sets from the
equal counts in \textup{(D.45.31)} retains equal
counts with \(\varphi(K_j)=S\), over arbitrary
finite bases.

\paragraph{The full normalizer count and generation.}
For a surface map on \(H_{j,D_*}\) with
\(\varphi(K_j)=S\), put
\(N_\varphi=\ker(\varphi|_{K_j})\).
Its normalizer in \(\widehat F\) contains
\(H_{j,D_*}\), and contains \(K_j\); it is
therefore one of the actual \(H_{j,C}\).
It contains \(H_{j,C}\) exactly when \(\varphi\)
has a unique extension to a map
\(H_{j,C}\to\operatorname{Aut}(S)\).
To prove this, if \(H_{j,C}\) normalizes the
kernel, conjugation on \(K_j/N_\varphi\simeq S\)
defines that extension.
It agrees with \(\varphi\) on \(H_{j,D_*}\)
because \(C_{\operatorname{Aut}(S)}(S)=1\).
Its outer coordinate factors through \(C\)
and is a unique \(\epsilon_C\) extending
\(\epsilon_*\). It is onto
\(S\rtimes\operatorname{im}(\epsilon_C)\),
since its \(K_j\)-image is \(S\).
Conversely an extension normalizes the kernel.
The same centralizer equality proves uniqueness.

Let \(M_{j,C}\) be the sum over all the
\(\epsilon_C\) in \textup{(D.45.24)} of the
retained counts after \textup{(D.45.32)}.
The previous paragraphs prove
\(M_{0,C}=M_{1,C}\) for every actual subgroup
\(C\geq D_*\).
The injections \(j_C\) and their base-coordinate
comparison prove that \(M_{j,C}\) counts precisely
those original surface maps whose kernel
normalizer contains \(H_{j,C}\).
Different extending characters are disjoint
because the extending map is unique.
Finite subgroup-poset inversion now gives the
exact number with minimal normalizer:
\[
        e_j=\sum_{C\geq D_*}\mu(D_*,C)\,M_{j,C},
                    \qquad e_0=e_1.
 \tag{D.45.35}
\]
Every summand refers to the displayed surface,
embedding and boundary maps; independent
isomorphisms of those surfaces suffice for this
cardinality equality.

The original epimorphism \(q_E\) has
\(q_E(K_0)=A\). Its coordinate kernels on
\(K_0\) are pairwise distinct and have precisely
the stated factor stabilizers, so \(e_0>0\).
For a map counted by \(e_1\), its conjugate
\(K_1\)-kernels at distinct factors are pairwise
distinct, by the exact normalizer condition.
Its \(K_1\)-image is then the full product
\(S^\Omega\): a proper subdirect product of
nonabelian simple factors identifies two factors
by an isomorphism and would identify those two
coordinate kernels. This follows inductively
from the two-factor Goursat lemma and the fact
that normal subgroups of a product of
nonabelian simple groups are products of factors.
The base image is \(D\), hence the resulting
map is onto \(E\).
Conversely every epimorphism to \(E\) has these
same full coordinate images and distinct
kernels.
The literal gauge count
\textup{(D.45.23)} therefore proves that the
two complete epimorphism fibres have the same
positive size
\[
                         |S|^{|\Omega|-1}e_j .
 \tag{D.45.36}
\]
This proves \textup{(D.45.4)}, including every
specified base conjugator.

\paragraph{The actual \(\operatorname{PSL}_2(32)\) source.}
We give a direct matrix proof of the group facts
used in \textup{(D.45.5)}.
Let \(k=\mathbf F_{32}\), and put \(S=\operatorname{SL}_2(k)\).
Its centre is trivial in characteristic two.
The square map of \(k^\times\) is bijective, so scaling
each invertible two-by-two matrix to determinant one
identifies \(S\) with \(\operatorname{PGL}_2(k)\).
Its action on \(\mathbf P^1(k)\) is sharply
three-transitive: the unique fractional linear map
between two ordered triples of distinct points has
such a determinant-one representative.
In particular it is faithful and primitive.

The upper unipotent subgroup
\(U=\{u(b):b\in k\}\) is abelian and normal in
the point stabilizer at infinity.
Its conjugates generate \(S\), by upper and lower
matrix elimination.
For \(h(a)=\operatorname{diag}(a,a^{-1})\),
\[
                 [h(a),u(b)]=u((a^2-1)b).
 \tag{D.45.37}
\]
Choose \(a\ne1\). Every element of \(U\) is then
a commutator; its conjugates prove that \(S\)
is perfect.
For a nontrivial normal subgroup \(M\), primitivity
and faithfulness make \(M\) transitive.
Thus every \(g\in S\) can be written \(g=mh\),
with \(m\in M\) and \(h\) in the stabilizer of
infinity.
Modulo \(M\), every conjugate of \(U\) therefore
has image equal to that same abelian image of
\(U\). These conjugates generate \(S\), so
\(S/M\) is abelian; perfectness gives \(M=S\).
This proves simplicity with the actual action.

The group \(U\) has order \(32\), the full
two-part of \(|S|=32(32^2-1)\).
Its normalizer is the upper triangular point
stabilizer of order \(32\cdot31\).
The \(33\) Sylow-two subgroups are therefore
exactly the conjugates indexed by
\(\mathbf P^1(k)\), each with its unique fixed
point.
An automorphism of \(S\) permutes these actual
subgroups and gives a permutation normalizing
the faithful projective action.
This map on automorphisms is injective:
the action of the image of any group element
is determined by its conjugation permutation
on the Sylow subgroups.
Conversely a permutation normalizing \(S\)
induces an automorphism of this faithful
permutation group.

Compose such a permutation by the unique
element of \(S\) making it fix infinity, zero
and one. It normalizes \(U\), the translation
group on \(k\), so the resulting permutation
\(\eta:k\to k\) satisfies
\[
                \eta(a+b)=\eta(a)+\eta(b).
 \tag{D.45.38}
\]
It also normalizes the stabilizer of infinity
and zero, the full group of nonzero scalar
multiplications. Since \(\eta(1)=1\), this
gives
\[
                \eta(ab)=\eta(a)\eta(b).
 \tag{D.45.39}
\]
Thus it is an actual field automorphism.
The field automorphisms of \(k/\mathbf F_2\)
are the five powers of squaring: the image of
an element of degree five has at most five
choices, and these five powers are distinct.
This proves the full equality
\(\operatorname{Aut}(S)=S\rtimes C_5\) and the
literal matrix generator \(t\).
It also proves the split outer quotient and
its order, with no uncomputed automorphism
enumeration.
The same identification is stated in the
opened primary account
\cite[\S3, preceding Lemma~3.1; Lemma~3.2]{ACyclicOuterWhite2011};
the preceding argument supplies its exact
group and field maps for the present source.

\paragraph{The fixed pro-five target and coupled generation.}
Take the actual D.37 quotient
\(\rho_n:\widehat F\twoheadrightarrow P_n\).
For \(q_1=\beta_Q\rho_n\), its three base
conjugators are exactly \(1,f_Q,h_Q\) of
D.37.14--19. Applying the proved theorem to
any original epimorphism above \(\rho_n\)
retains these specified values in
\textup{(D.45.4)}.
Every homomorphism from \(S^\Omega\) to a
finite five-group is trivial: its image is
perfect, while a finite five-group is solvable.
Hence every pro-five quotient of \(E\)
factors through its displayed \(P_n\), which
is itself a five-group.
This proves the exact maximal pro-five
comparison, and retains all the D.37
parameters and all their powers.

More generally let \(L\twoheadrightarrow D\)
have solvable kernel, and fix an epimorphism
of the same source into \(L\) with its three
peripheral conjugators above the \(h_i\).
Pair that map with \textup{(D.45.4)}.
The conjugating elements pair in the actual
fibre product \(E\times_D L\).
The two projections are onto. Their possible
missing common kernel quotient would be a
quotient of \(S^\Omega\), hence perfect, and
also a quotient of the solvable kernel of
\(L\to D\). It is therefore trivial.
Goursat proves that the paired map is onto
the full actual fibre product.
\end{proof}

The prime order in \textup{(D.45.1)} is used in
the independent-phase maps: every nonzero
character value generates the outer image,
and the exponent \(e\) in
\textup{(D.45.15)} is invertible.
The remaining extension, Shapiro, subtraction
and normalizer arguments use the stated
split automorphism group and simplicity.
The \(\operatorname{PSL}_2(32)\) calculation
supplies these hypotheses with the literal
field automorphism. It does not assign
outer coordinates by class preservation of
\(\beta_Q\).

The certificate \texttt{a\_cyclic\_outer\_phase.py}
retains every free-word map and its inverse,
all \(24\) nonzero handle characters over
\(\mathbf F_5\), and an exact model of all
\(32736\) matrices in \(S\).
It checks the sharp three-transitive action,
all \(33\) \(S\)-classes, the six field orbits
of length five, and the three \(S\)-classes
in each of the four nonidentity outer
cosets. It also gives ten actual surjective
surface maps whose selected boundary
traverses all five field phases, with the
inverse source words retained.
The companion
\texttt{a\_cyclic\_outer\_shapiro.py}
uses the actual order-\(125\) Heisenberg base,
its central order-five factor stabilizer,
and every one of the five character
extensions to an order-\(25\) overgroup.
It verifies both signs in
\textup{(D.45.19)}, all intermediate
embeddings with actual
\(\operatorname{Aut}(S)\) multiplication,
and all the chosen base-conjugator
comparisons. Its central character does
not extend to the whole Heisenberg group,
as witnessed by the displayed commutator;
that nonextension is retained in the
intermediate count.
